\begin{table}[htbp]
\centering
\caption{各類型身心狀態平均數與調整後組間差異}
\label{tab:main}
\footnotesize
\setlength{\tabcolsep}{4pt}
\renewcommand{\arraystretch}{1.12}
\begin{tabular}{l rrrr r rrrr}
\toprule
& \multicolumn{4}{c}{各類型平均數（標準差）} & & \multicolumn{4}{c}{調整後組間差異（四象限分類模型）} \\
\cmidrule(lr){2-5} \cmidrule(lr){7-10}
身心狀態指標 & 雙高 & 線上優勢 & 線下優勢 & 雙低 & $\eta^2$ & 線上－雙低 & 線上－線下 & 雙高－線上 & 雙高－線下 \\
\midrule
幸福感 & 2.78 & 1.97 & 2.92 & 2.18 & .096 & $-0.17^{***}$ & $-0.89^{***}$ & $0.78^{***}$ & $-0.12^{**}$ \\
 & (1.17) & (1.06) & (1.07) & (1.25) & & (0.040) & (0.040) & (0.034) & (0.035) \\[2pt]
自尊 & 3.26 & 2.52 & 3.32 & 2.71 & .170 & $-0.14^{***}$ & $-0.75^{***}$ & $0.68^{***}$ & $-0.07^{**}$ \\
 & (0.68) & (0.78) & (0.62) & (0.83) & & (0.027) & (0.026) & (0.023) & (0.021) \\[2pt]
生活滿意度 & 3.28 & 2.74 & 3.37 & 2.92 & .099 & $-0.17^{***}$ & $-0.62^{***}$ & $0.51^{***}$ & $-0.11^{***}$ \\
 & (0.70) & (0.77) & (0.62) & (0.81) & & (0.027) & (0.025) & (0.023) & (0.020) \\[2pt]
好友支持 & 3.23 & 3.02 & 3.20 & 2.74 & .072 & $0.20^{***}$ & $-0.21^{***}$ & $0.27^{***}$ & $0.06^{**}$ \\
 & (0.63) & (0.69) & (0.65) & (0.83) & & (0.026) & (0.024) & (0.021) & (0.020) \\[2pt]
孤獨感$^{\dagger}$ & 0.43 & 0.91 & 0.26 & 0.55 & .044 & $0.33^{***}$ & $0.59^{***}$ & $-0.43^{***}$ & $0.16^{***}$ \\
 & (0.93) & (1.24) & (0.69) & (1.07) & & (0.040) & (0.036) & (0.035) & (0.025) \\[2pt]
心理困擾$^{\dagger}$ & 0.50 & 0.97 & 0.36 & 0.63 & .063 & $0.29^{***}$ & $0.55^{***}$ & $-0.42^{***}$ & $0.14^{***}$ \\
 & (0.76) & (0.93) & (0.51) & (0.84) & & (0.030) & (0.027) & (0.026) & (0.019) \\[2pt]
\bottomrule
\end{tabular}

\vspace{0.4em}
\parbox{\linewidth}{\scriptsize 註：平均數為原始分數，六項指標的 ANOVA 與 Kruskal-Wallis 檢定經 FDR 校正後皆 $p<.001$。組間差異為納入背景控制變數之四象限分類模型的係數線性組合，括號內為 HC3 穩健標準誤，$p$ 值經 FDR 校正；「線上」「線下」分別指線上優勢型與線下優勢型。$^{*}p<.05$，$^{**}p<.01$，$^{***}p<.001$。$^{\dagger}$分數越高代表狀態越差。}
\end{table}
